\chapter{Martyngały: rachunek, strategie i nierówności}
\label{ch:martyngaly}
\paragraph{Odsyłacze międzytomowe.} Zastosowania martyngałów w diagnostyce MCMC omawia tom IV, rozdział~\ref{IV-ch:mcmc}, a martyngały ciągłe tom II, rozdział~\ref{II-ch:mart-cont}.


\section{Średni przyrost równy zero nie oznacza braku zmian}

Martyngał może bardzo silnie fluktuować: nie jest
procesem stałym i kolejne wartości nie muszą być
niezależne. Warunek
$\E[M_{n+1}-M_n\mid\mathcal F_n]=0$
mówi jedynie, że \emph{po uwzględnieniu całej
aktualnie dostępnej informacji} nie ma
przewidywalnego zysku ani straty.
Z definicji martyngału (rozdział~\ref{ch:filtracje})
oraz własności wieży (rozdział~\ref{ch:condexp})
otrzymujemy $\E M_{n+1}=\E M_n$,
a przez indukcję $\E M_n=\E M_0$ dla każdego $n$,
ale sama stałość średniej nie wystarcza do
martyngałowości.

\begin{przyklad}[Stała średnia bez martyngału]
Niech $\xi$ przyjmuje $\pm1$ z prawdopodobieństwem
$1/2$ i połóżmy $X_0=0$, $X_1=\xi$,
$X_2=-\xi$, $X_n=X_2$ dla $n\geq2$.
Mamy $\E X_n=0$ dla każdego $n$, niezależnie
od wyboru filtracji. W naturalnej filtracji
$\mathcal F_n=\sigma(X_0,\ldots,X_n)$ wszystkie
$X_n$ są adaptowane i ograniczone. Jednak
$\mathcal F_1=\sigma(\xi)$, więc
$\E[X_2\mid\mathcal F_1]=-\xi\ne\xi=X_1$ p.n.
Średnia bezwarunkowa gubi informację o znaku
pierwszego wyniku.
\end{przyklad}

\begin{przyklad}[Martyngał dla dowolnej zmiennej końcowej]
Niech $Y\in L^1(\Prob)$ i niech $(\mathcal F_n)$
będzie filtracją. Wówczas
$M_n=\E[Y\mid\mathcal F_n]$ (wybieramy wersje
$\mathcal F_n$-mierzalne) jest martyngałem:
$\E|M_n|\leq\E|Y|$ z warunkowej nierówności
Jensena, a z własności wieży
$\E[M_{n+1}\mid\mathcal F_n]=
\E[\E[Y\mid\mathcal F_{n+1}]\mid\mathcal F_n]=M_n$ p.n.
Interpretujemy $M_n$ jako aktualną prognozę
przyszłej wielkości $Y$ na podstawie informacji
do chwili $n$.
\end{przyklad}

\begin{przyklad}[Losowe mnożenie]
Niech $Z_1,Z_2,\ldots$ będą niezależne, nieujemne,
$\E Z_k=1$, i zdefiniujmy $M_0=1$,
$M_n=\prod_{k=1}^n Z_k$.
Przyjmijmy $\mathcal F_n=\sigma(Z_1,\ldots,Z_n)$,
a $\mathcal F_0=\{\varnothing,\Omega\}$.
Proces jest adaptowany, a z niezależności
i nieujemności czynników wynika
$\E M_n=\prod_{k=1}^n\E Z_k=1$.
Ponieważ $Z_{n+1}$ jest niezależna od $\mathcal F_n$,
reguła wyciągania mierzalnego czynnika daje
$\E[M_{n+1}\mid\mathcal F_n]
=M_n\E Z_{n+1}=M_n$ p.n.
Niekoniecznie $M_n\to1$. Na przykład gdy
każde $Z_k$ przyjmuje $0$ i $2$ z prawdopodobieństwem
$1/2$, mamy $\Prob(M_n\ne0)=2^{-n}$.
Zdarzenia $\{M_n\ne0\}$ są malejące, więc
z ciągłości miary od góry wynika, że p.n.
któryś czynnik jest zerowy i od tej chwili
$M_n=0$. Mamy zatem $M_n\to0$ p.n.,
mimo $\E M_n=1$ dla każdego $n$. Do tej
subtelności wrócimy w rozdziale~\ref{ch:nadmartyngaly}.
\end{przyklad}

\begin{przyklad}[Martyngał wykładniczy]
Niech $\xi_1,\xi_2,\ldots$ będą niezależne i jednakowo
rozłożone, a dla ustalonego $t\in\R$
niech $m(t)=\E e^{t\xi_1}<\infty$.
Ponieważ $\xi_1$ jest rzeczywista, $e^{t\xi_1}>0$,
więc $m(t)>0$. Dla $S_0=0$,
$S_n=\sum_{k=1}^n\xi_k$ proces
\[
M_n(t)=\frac{e^{tS_n}}{m(t)^n},\qquad M_0(t)=1,
\]
jest nieujemnym martyngałem względem filtracji
$\mathcal F_n=\sigma(\xi_1,\ldots,\xi_n)$
(z trywialnym $\mathcal F_0$), która jest także
naturalną filtracją sum $(S_n)$.
Z niezależności wynika $\E M_n(t)=1$, zatem
proces jest całkowalny, a z konstrukcji adaptowany.
Wyciągając mierzalny czynnik z warunkowej
wartości oczekiwanej i korzystając z niezależności,
otrzymujemy
\[
\E[M_{n+1}(t)\mid\mathcal F_n]
=M_n(t)\frac{\E e^{t\xi_{n+1}}}{m(t)}
=M_n(t)\quad\text{p.n.}
\]
Dla spaceru symetrycznego
$m(t)=(e^t+e^{-t})/2=\cosh t$.
Takie martyngały umożliwiają szacowanie
prawdopodobieństw trafienia w odległy próg.
\end{przyklad}

\section{Różnice martyngałowe i ortogonalność}

\begin{definicja}[Różnice martyngałowe]
Ciąg $(D_n)_{n\geq1}$ rzeczywistych zmiennych
losowych na $(\Omega,\mathcal F,\Prob)$ nazywamy
\emph{ciągiem różnic martyngałowych} względem
filtracji $(\mathcal F_n)$, jeśli dla każdego $n\geq1$
zmienna $D_n$ jest całkowalna, $\mathcal F_n$-mierzalna
oraz $\E[D_n\mid\mathcal F_{n-1}]=0$ p.n.
\end{definicja}

\begin{stwierdzenie}[Reprezentacja przez różnice martyngałowe]
\label{prop:rp1-14-differences}
Jeśli $(M_n)$ jest martyngałem względem
$(\mathcal F_n)$, to $D_n=M_n-M_{n-1}$ tworzą
ciąg różnic martyngałowych.
Odwrotnie, dla dowolnego ciągu różnic
martyngałowych $(D_n)$ oraz całkowalnej
$\mathcal F_0$-mierzalnej zmiennej $M_0$,
proces $M_n=M_0+\sum_{k=1}^nD_k$ jest
martyngałem względem tej samej filtracji.
\end{stwierdzenie}
\begin{proof}
Dla przyrostów martyngału całkowalność wynika
z nierówności trójkąta, a mierzalność z adaptacji
i inkluzji $\mathcal F_{n-1}\subseteq\mathcal F_n$.
Liniowość warunkowej wartości oczekiwanej daje
$\E[D_n\mid\mathcal F_{n-1}]=M_{n-1}-M_{n-1}=0$ p.n.
Odwrotnie, skończona suma całkowalnych zmiennych
jest całkowalna i $\mathcal F_n$-mierzalna,
a $\E[M_n\mid\mathcal F_{n-1}]
=M_{n-1}+\E[D_n\mid\mathcal F_{n-1}]=M_{n-1}$ p.n.
\end{proof}

\begin{twierdzenie}[Ortogonalność przyrostów w $L^2$]
Niech $(M_n)$ będzie martyngałem takim, że
$M_n\in L^2$ dla wszystkich $n$. Dla
$D_k=M_k-M_{k-1}$, $k\geq1$, mamy
$D_k\in L^2$ oraz $\E[D_iD_j]=0$ dla
różnych dodatnich indeksów $i,j$. Ponadto
\[
\E[(M_n-M_0)^2]=\sum_{k=1}^n\E D_k^2,
\qquad
\E M_n^2=\E M_0^2+\sum_{k=1}^n\E D_k^2.
\]
\end{twierdzenie}
\begin{proof}
Każdy przyrost należy do $L^2$, ponieważ
jest różnicą dwóch elementów tej przestrzeni.
Załóżmy $i<j$. Zmienna $D_i$ jest
$\mathcal F_{j-1}$-mierzalna na mocy
$\mathcal F_i\subseteq\mathcal F_{j-1}$,
a iloczyn $D_iD_j$ jest całkowalny na mocy
nierówności Cauchy'ego--Schwarza. Pełna wersja reguły
wyciągania czynnika z warunkowania
(rozdział~\ref{ch:condexp}) daje
$\E[D_iD_j]=\E[D_i\E[D_j\mid\mathcal F_{j-1}]]=0$.
Podobnie $\E[M_0D_j]=0$, gdyż $M_0$ jest
$\mathcal F_{j-1}$-mierzalna, a $M_0D_j\in L^1$
z nierówności Cauchy'ego--Schwarza.
Rozwijając kwadrat sumy $\sum_{k=1}^nD_k$
oraz osobno kwadrat $M_0+\sum_{k=1}^nD_k$,
otrzymujemy obie tożsamości, ponieważ wszystkie
wyrazy mieszane mają średnią zero.
\end{proof}

\begin{uwaga}
Dla wycentrowanych zmiennych w $L^2$ niezależność
implikuje ortogonalność, ponieważ wartość oczekiwana
iloczynu jest iloczynem wartości oczekiwanych.
Odwrotna implikacja nie zachodzi. Różnice
martyngałowe w $L^2$ mają średnie zero, więc
ortogonalność oznacza tutaj także zerowe kowariancje,
ale nie wyklucza zależności od wcześniejszych obserwacji.
\end{uwaga}

\begin{przyklad}[Zależne, lecz ortogonalne przyrosty]
Niech $\xi_k$ będą niezależnymi symetrycznymi
znakami, $D_1=\xi_1$, $D_2=\xi_1\xi_2$,
$D_3=\xi_1(1+\xi_2)\xi_3$ i
$D_n=\xi_n$ dla $n\geq4$.
Względem $\mathcal F_n=\sigma(\xi_1,\ldots,\xi_n)$
każdy przyrost ma warunkową średnią zero,
bo jest iloczynem ograniczonej zmiennej
$\mathcal F_{n-1}$-mierzalnej i niezależnej od
$\mathcal F_{n-1}$ zmiennej $\xi_n$ o średniej zero.
Wszystkie przyrosty są ograniczone, a więc
należą do $L^2$. Mamy
$\{|D_3|=0\}=\{\xi_2=-1\}=\{D_1D_2=-1\}$.
Zdarzenie to ma prawdopodobieństwo $1/2$ i należy
zarówno do $\sigma(D_3)$, jak i do $\sigma(D_1,D_2)$.
Gdyby $D_3$ była niezależna od pary $(D_1,D_2)$,
jego prawdopodobieństwo musiałoby być równe
własnemu kwadratowi, co jest niemożliwe.
Ciąg przyrostów nie jest więc wspólnie niezależny.
Mimo to wszystkie różne przyrosty są ortogonalne
na mocy udowodnionego twierdzenia.
\end{przyklad}

\section{Suma przewidywalnie ważonych przyrostów}

\begin{twierdzenie}[Transformacja martyngałowa]
Niech $(M_n)$ będzie rzeczywistym martyngałem
względem $(\mathcal F_n)$ i $(H_n)_{n\geq1}$
rzeczywistym procesem przewidywalnym względem
tej samej filtracji.
Załóżmy, że każda zmienna
$H_n(M_n-M_{n-1})$ jest całkowalna
(np. wszystkie $H_n$ są ograniczone).
Wtedy
\[
(H\cdot M)_0=0,\qquad
(H\cdot M)_n=\sum_{k=1}^nH_k(M_k-M_{k-1})
\]
jest martyngałem względem $(\mathcal F_n)$.
Jeśli zamiast martyngału przyjmiemy podmartyngał
(odpowiednio nadmartyngał), zachowamy założenie
całkowalności iloczynów i założymy dodatkowo
$H_n\geq0$ p.n. dla każdego $n$, to transformacja
jest podmartyngałem (odpowiednio nadmartyngałem).
\end{twierdzenie}
\begin{proof}
Każdy składnik sumy jest całkowalny na mocy
założenia oraz $\mathcal F_k$-mierzalny, ponieważ
$H_k$ jest $\mathcal F_{k-1}$-mierzalna.
Skończona suma jest więc całkowalna
i $\mathcal F_n$-mierzalna.
Ponieważ $H_{n+1}$ jest $\mathcal F_n$-mierzalna,
reguła wyciągania czynnika, w jej wersji
z warunkiem całkowalności iloczynu, daje
\[
\E[(H\cdot M)_{n+1}-(H\cdot M)_n
\mid\mathcal F_n]
=H_{n+1}\E[M_{n+1}-M_n\mid\mathcal F_n].
\]
W martyngale prawa strona jest zerowa; w
pod- i nadmartyngale przy $H_{n+1}\geq0$
ma odpowiedni znak.
\end{proof}

\begin{przyklad}[Zmienianie stawek nie tworzy przewidywalnego zysku]
Niech $M_n$ oznacza skumulowany wynik uczciwej gry,
a $H_n$ stawkę wybraną na podstawie historii
$\mathcal F_{n-1}$. Jeżeli ryzyko jest na tyle
kontrolowane, że $H_n\Delta M_n$ jest całkowalny,
proces sumy zysków jest martyngałem.
\emph{Nie} wynika z tego, że wynik gry nie może
być dodatni na określonej trajektorii, tylko że
nie ma dodatniego \emph{warunkowego oczekiwanego}
przyrostu. Twierdzenie dotyczy każdej skończonej chwili $n$
przy podanej całkowalności iloczynów.
Wartość transformacji w nieograniczonym czasie
zatrzymania wymaga dodatkowych założeń
(rozdział~\ref{ch:filtracje}). Nie wolno więc
wnioskować o jej średniej wyłącznie na podstawie
martyngałowości w skończonych chwilach.
\end{przyklad}

\section{Rozkład Dooba i energia martyngału}

\begin{twierdzenie}[Rozkład Dooba w czasie dyskretnym]
Niech $(X_n)$ będzie rzeczywistym procesem
adaptowanym względem $(\mathcal F_n)$ i niech
$X_n\in L^1$ dla każdego $n$.
Wtedy istnieje rozkład $X_n=M_n+A_n$,
gdzie $M$ jest martyngałem względem $(\mathcal F_n)$,
$A_0=0$, każda zmienna $A_n$ jest całkowalna,
a $(A_n)_{n\geq1}$ jest przewidywalny.
Rozkład jest jednoznaczny z dokładnością
do nierozróżnialności procesów, tj. oba rozkłady
pokrywają się we wszystkich chwilach na jednym
zbiorze pełnej miary.
Jeśli $X$ jest podmartyngałem, to $A$ jest
niemalejący p.n.; jeśli $X$ jest nadmartyngałem,
to $A$ jest nierosnący p.n.
Proces $A$ nazywamy przewidywalną częścią
rozkładu Dooba.
\end{twierdzenie}
\begin{proof}
Definiujemy
\[
A_0=0,\qquad
A_n=\sum_{k=1}^n\E[X_k-X_{k-1}\mid\mathcal F_{k-1}],
\qquad M_n=X_n-A_n.
\]
Wybieramy całkowalne $\mathcal F_{k-1}$-mierzalne
wersje warunkowych wartości oczekiwanych.
Dla $k\leq n$ mamy $\mathcal F_{k-1}\subseteq
\mathcal F_{n-1}$, więc $A_n$ jest całkowalna
i $\mathcal F_{n-1}$-mierzalna dla $n\geq1$.
Jeżeli $X$ jest podmartyngałem, przyrosty
$\Delta A_k=A_k-A_{k-1}$ są nieujemne p.n.;
dla nadmartyngału są niedodatnie p.n.
Przeliczalność indeksów pozwala wybrać jeden
zbiór pełnej miary, na którym wszystkie te
nierówności zachodzą jednocześnie.
Proces $M$ jest całkowalny i adaptowany, a
$\E[M_n-M_{n-1}\mid\mathcal F_{n-1}]
=\E[X_n-X_{n-1}\mid\mathcal F_{n-1}]
-(A_n-A_{n-1})=0$.
Jeśli $X=M'+A'$ byłby drugim rozkładem,
to z $\Delta A'_n=\Delta X_n-\Delta M'_n$
po warunkowaniu otrzymujemy
$\Delta A'_n=\E[\Delta X_n\mid\mathcal F_{n-1}]$.
Wraz z $A'_0=0$ daje to $A'_n=A_n$ p.n.
dla każdego $n$, a stąd $M'_n=M_n$ p.n.
Przecinając przeliczalną rodzinę zbiorów
pełnej miary, otrzymujemy nierozróżnialność.
\end{proof}

\begin{definicja}[Wariacja kwadratowa i przewidywalna wariacja kwadratowa]
Dla rzeczywistego martyngału $(M_n)$ takiego,
że $M_n\in L^2$ dla każdego $n$, oznaczmy
$\Delta M_k=M_k-M_{k-1}$.
Definiujemy $[M]_0=\langle M\rangle_0=0$ oraz
\[
[M]_n=\sum_{k=1}^n(\Delta M_k)^2,\qquad
\langle M\rangle_n=\sum_{k=1}^n
\E[(\Delta M_k)^2\mid\mathcal F_{k-1}].
\]
Procesy te nazywamy odpowiednio \emph{wariacją
kwadratową} i \emph{przewidywalną wariacją
kwadratową} martyngału $M$.
\end{definicja}

\begin{stwierdzenie}[Kompensacja wariacji kwadratowej]
Procesy $(M_n^2-\langle M\rangle_n)$ oraz
$([M]_n-\langle M\rangle_n)$ są martyngałami.
Ponadto dla każdego $n$
\[
\E\bigl[[M]_n\bigr]=\E\langle M\rangle_n
=\E M_n^2-\E M_0^2.
\]
\end{stwierdzenie}
\begin{proof}
Obie wariacje są całkowalne i adaptowane,
ponieważ $\Delta M_k\in L^2$.
Warunkowa nierówność Jensena daje
$\E[M_n^2\mid\mathcal F_{n-1}]
\geq M_{n-1}^2$ p.n., więc $(M_n^2)$ jest
podmartyngałem. Z rozwinięcia kwadratu
i reguły wyciągania mierzalnego czynnika
(uzasadnionej całkowalnością $M_{n-1}\Delta M_n$
przez nierówność Cauchy'ego--Schwarza)
wynika
\[
\E[M_n^2-M_{n-1}^2\mid\mathcal F_{n-1}]
=\E[(\Delta M_n)^2\mid\mathcal F_{n-1}]
\quad\text{p.n.}
\]
Z rozkładu Dooba częścią przewidywalną $(M_n^2)$
jest zatem $\langle M\rangle$, co dowodzi
pierwszej własności martyngałowej.
Dla drugiej, przyrost procesu $[M]-\langle M\rangle$
jest równy
$(\Delta M_n)^2-\E[(\Delta M_n)^2\mid\mathcal F_{n-1}]$,
czyli ma warunkową średnią zero.
Biorąc wartości oczekiwane obu martyngałów
i korzystając z ich wartości w chwili zero,
otrzymujemy podane tożsamości.
\end{proof}

Wariacja kwadratowa rejestruje rzeczywiste
kwadraty przyrostów, a przewidywalna wariacja
kwadratowa sumuje ich warunkowe drugie momenty.
Nie muszą być równe ścieżkowo, mimo równych średnich.

\begin{przyklad}[Zależna zmienność]
Niech $\xi_k\in\{-1,1\}$ będą niezależne
i symetryczne, $\mathcal F_n=\sigma(\xi_1,\ldots,\xi_n)$
(z trywialnym $\mathcal F_0$), a każda zmienna
$\sigma_k>0$ niech będzie $\mathcal F_{k-1}$-mierzalna
oraz spełnia $\E\sigma_k^2<\infty$.
Przyjmijmy $M_0=0$ i $M_n=\sum_{k=1}^n\sigma_k\xi_k$.
Przyrosty należą do $L^2$, ponieważ
$(\sigma_k\xi_k)^2=\sigma_k^2$. Ich warunkowe
średnie wynoszą $\sigma_k\E\xi_k=0$ p.n.
Ze stwierdzenia~\ref{prop:rp1-14-differences}
wynika, że $M$ jest martyngałem w $L^2$.
Ponieważ kwadrat przyrostu jest już
$\mathcal F_{k-1}$-mierzalny, w tym szczególnym
przypadku $[M]_n=\langle M\rangle_n=\sum_{k=1}^n\sigma_k^2$
p.n.
Zmienność przyrostu może zależeć od historii;
warunek martyngałowy wymaga tylko zaniku
warunkowej średniej.
\end{przyklad}

\section{Nierówności maksymalne Dooba}

\begin{twierdzenie}[Nierówność maksymalna $L^1$]
Niech $(X_k)_{0\leq k\leq n}$ będzie nieujemnym
podmartyngałem względem $(\mathcal F_k)$.
Dla każdego $\lambda>0$
\[
\lambda\Prob\left(\max_{0\leq k\leq n}X_k\geq\lambda\right)
\leq\E\left[X_n\chi_{\{\max_{0\leq k\leq n}X_k\geq\lambda\}}\right]
\leq\E X_n.
\]
W szczególności dla martyngału $(M_k)$
$\Prob(\max_{k\leq n}|M_k|\geq\lambda)
\leq\E|M_n|/\lambda$.
\end{twierdzenie}
\begin{proof}
Niech $A=\{\max_{0\leq k\leq n}X_k\geq\lambda\}$
i $\tau=\inf\{k\leq n:X_k\geq\lambda\}$,
z umową $\tau=n$ na ścieżkach bez przekroczenia.
Rozbijmy $A$ na rozłączne zdarzenia
$A_k=\{X_0<\lambda,\ldots,X_{k-1}<\lambda,
X_k\geq\lambda\}$ dla $k\leq n$.
Mamy $A_k\in\mathcal F_k$ i
$X_k\geq\lambda$ na $A_k$.
Z podmartyngałowości i wieży
$\E[X_n\chi_{A_k}]\geq\E[X_k\chi_{A_k}]
\geq\lambda\Prob(A_k)$.
Sumujemy po $k=0,\ldots,n$, dostając pierwszą
nierówność. Druga wynika z nieujemności $X_n$.
Dla martyngału $|M_k|$ jest podmartyngałem
przez warunkową nierówność Jensena.
\end{proof}

\begin{twierdzenie}[Nierówność maksymalna $L^p$]
Niech $1<p<\infty$, $n\in\N_0$ i
$(X_k)_{k\leq n}$ będzie nieujemnym
podmartyngałem z $X_n\in L^p$. Wtedy
\[
\left\|\max_{0\leq k\leq n}X_k\right\|_p
\leq\frac{p}{p-1}\,\|X_n\|_p.
\]
Dla martyngału $M$ z $M_n\in L^p$ można
zastąpić $X_k$ przez $|M_k|$.
\end{twierdzenie}
\begin{proof}
Z nierówności podmartyngałowej $X_k\leq
\E[X_n\mid\mathcal F_k]$ oraz warunkowego
Jensena wynika najpierw, że
$\E X_k^p\leq\E X_n^p$ dla wszystkich $k\leq n$;
w szczególności każde $X_k$ ma skończony $p$-ty moment.
Oznaczmy $Y=\max_{k\leq n}X_k$.
Z poprzedniego twierdzenia, dla każdego $t>0$,
$t\Prob(Y\geq t)\leq\E[X_n\chi_{\{Y\geq t\}}]$.
Najpierw zastąpmy $Y$ przez $Y\wedge L$,
aby wszystkie całki były skończone.
Wzór warstwowy i twierdzenie Tonellego dają
\begin{align*}
\E[(Y\wedge L)^p]
&=p\int_0^L t^{p-1}\Prob(Y\geq t)\,\dd\lambda_1(t)\\
&\leq p\int_0^L t^{p-2}
\E[X_n\chi_{\{Y\geq t\}}] \,\dd\lambda_1(t)\\
&=\frac{p}{p-1}\E\bigl[X_n(Y\wedge L)^{p-1}\bigr]\\
&\leq\frac{p}{p-1}\|X_n\|_p
\|Y\wedge L\|_p^{p-1},
\end{align*}
gdzie w ostatnim kroku zastosowano nierówność
Höldera. Jeśli $\|Y\wedge L\|_p=0$, teza jest
oczywista; w innym przypadku dzielimy przez
$\|Y\wedge L\|_p^{p-1}$, a następnie
przechodzimy z $L\to\infty$ przez monotoniczną
zbieżność.
\end{proof}

\begin{przyklad}[Kontrola maksimum trajektorii]
Jeżeli $S_n$ jest symetrycznym spacerem losowym,
to $\E S_n^2=n$. Z nierówności Dooba dla $p=2$
otrzymujemy
$\E(\max_{k\leq n}|S_k|)^2\leq4n$.
Nierówność Markowa daje
$\Prob(\max_{k\leq n}|S_k|\geq u)\leq4n/u^2$.
Kontrolujemy w ten sposób \emph{całą drogę},
a nie tylko wartość końcową $S_n$.
\end{przyklad}

\section{Nierówność Azumy--Hoeffdinga}

\begin{twierdzenie}[Koncentracja martyngału o ograniczonych przyrostach]
Niech $(M_k)_{0\leq k\leq n}$ będzie rzeczywistym
martyngałem. Załóżmy, że istnieją stałe
$c_k\geq0$ spełniające
$|M_k-M_{k-1}|\leq c_k$ p.n. dla $1\leq k\leq n$.
Wtedy dla każdego $u>0$ i
$V=\sum_{k=1}^n c_k^2>0$ zachodzi
\[
\Prob(M_n-M_0\geq u)\leq e^{-u^2/(2V)},
\quad
\Prob(|M_n-M_0|\geq u)\leq2e^{-u^2/(2V)}.
\]
Gdy $V=0$, przyrosty są zerowe p.n. i zdarzenie
odchylenia o dodatnie $u$ ma prawdopodobieństwo zero.
\end{twierdzenie}
\begin{proof}
Dla $D_k=\Delta M_k$ mamy
$\E[D_k\mid\mathcal F_{k-1}]=0$ oraz
$-c_k\leq D_k\leq c_k$.
Warunkowa wersja lematu Hoeffdinga dowodzona
jest dokładnie jak wersja bezwarunkowa, z
miarą prawdopodobieństwa zastąpioną warunkową
średnią i tym samym oszacowaniem wariancji
zmiennej na odcinku. Daje
\[
\E[e^{tD_k}\mid\mathcal F_{k-1}]
\leq e^{t^2c_k^2/2}\qquad(t\in\R).
\]
Można to również uzasadnić bez odwoływania się
do regularnych rozkładów warunkowych:
na przedziale $[-c_k,c_k]$ funkcja $e^{tx}$
leży poniżej cięciwy łączącej końce,
więc po warunkowaniu i użyciu zerowej średniej
otrzymuje się $\cosh(tc_k)\leq e^{t^2c_k^2/2}$.
Własność wieży daje rekurencyjnie
\[
\E e^{t(M_n-M_0)}\leq
\exp\left(\frac{t^2}{2}\sum_{k=1}^n c_k^2\right).
\]
Dla $t>0$ nierówność Markowa prowadzi do
$\Prob(M_n-M_0\geq u)\leq
\exp(-tu+t^2V/2)$.
Minimalizujemy przy $t=u/V$, otrzymując
$e^{-u^2/(2V)}$. Ten sam rachunek stosujemy do
$-(M_n-M_0)$ i dodajemy dwa ogony.
\end{proof}

\begin{przyklad}[Uczciwy spacer i prawdopodobieństwo odchylenia]
Dla niezależnych symetrycznych znaków
$S_n=\sum_{k\leq n}\xi_k$ i $c_k=1$.
Zatem dla $u>0$
$\Prob(|S_n|\geq u)\leq2e^{-u^2/(2n)}$.
Nierówność Czebyszewa dałaby jedynie
$\Prob(|S_n|\geq u)\leq n/u^2$.
Oszacowanie wykładnicze jest dużo mocniejsze
przy dużych odchyleniach, ale korzysta z ograniczenia
przyrostów, nie tylko z wariancji.
\end{przyklad}

\begin{przyklad}[Odkrywanie współrzędnych funkcji wielu losowań]
Niech $\xi_1,\ldots,\xi_n$ będą niezależne, a
$f(\xi_1,\ldots,\xi_n)$ całkowalna.
Definiujemy martyngał Dooba
$M_k=\E[f(\xi_1,\ldots,\xi_n)\mid
\sigma(\xi_1,\ldots,\xi_k)]$, gdzie $M_0=\E f$.
Jeżeli zmiana tylko $k$-tej współrzędnej
może zmienić $f$ o co najwyżej $c_k$,
to dla ustalonej historii $\xi_1,\ldots,\xi_{k-1}$
możliwe wartości $M_k$ leżą w przedziale długości
co najwyżej $c_k$. W konsekwencji
$|M_k-M_{k-1}|\leq c_k$; Azuma daje
\[
\Prob(|f-\E f|\geq u)\leq
2\exp\left(-\frac{u^2}{2\sum_kc_k^2}\right).
\]
Ostrzejszą stałą $2$ w wykładniku daje
warunkowy lemat Hoeffdinga zastosowany do
przyrostu $M_k$, którego \emph{przedział wartości}
ma długość $c_k$; wtedy otrzymuje się
$\Prob(|f-\E f|\geq u)\leq
2\exp(-2u^2/\sum_kc_k^2)$.
To znana nierówność ograniczonych różnic.
\end{przyklad}


\begin{uwaga}[Dlaczego w nierówności ograniczonych różnic jest lepsza stała?]
Samo oszacowanie $|\Delta M_k|\leq c_k$ mówi,
że przyrost leży w przedziale długości $2c_k$.
W przykładzie odkrywania współrzędnych
istnieje jednak przedział $[a_k,b_k]$
o długości $b_k-a_k\leq c_k$, wyznaczony przez
informację do chwili $k-1$, w którym leży
$\Delta M_k$. Dla zmiennej o warunkowej
średniej zero warunkowy lemat Hoeffdinga daje
\[
 \E[e^{t\Delta M_k}\mid\mathcal F_{k-1}]
 \leq\exp\bigl(t^2(b_k-a_k)^2/8\bigr)
 \leq\exp(t^2c_k^2/8).
\]
Po wymnożeniu oszacowań przez wieżę i minimalizacji
$-tu+t^2\sum_kc_k^2/8$ otrzymujemy
$\exp(-2u^2/\sum_kc_k^2)$ dla jednego ogona.
\end{uwaga}

\section{Proces Galtona--Watsona i znormalizowany martyngał}

\begin{przyklad}[Rozmnażanie i zmienna populacja]
Niech $Z_0=1$, a
$Z_{n+1}=\sum_{j=1}^{Z_n}\xi_{n,j}$, gdzie
$\xi_{n,j}$ są niezależnymi, jednakowo rozłożonymi
nieujemnymi całkowitymi zmiennymi losowymi,
niezależnymi od dotychczasowej historii,
o średniej $m=\E\xi_{n,j}>0$ i wariancji
$\sigma^2<\infty$. Filtrację $\mathcal F_n$
tworzą populacje i wszystkie losowania do
chwili $n-1$. Warunkowo przy $Z_n$ mamy
\[
\E[Z_{n+1}\mid\mathcal F_n]=mZ_n,
\qquad \Var(Z_{n+1}\mid\mathcal F_n)=\sigma^2 Z_n.
\]
Zatem $W_n=Z_n/m^n$ jest nieujemnym martyngałem.
Średnia $\E Z_n=m^n$ wynika indukcyjnie z wieży,
a więc $\E W_n=1$.

Jeśli $m>1$, to przez rozkład wariancji
\[
\E[(W_{n+1}-W_n)^2]
=\frac{\sigma^2\E Z_n}{m^{2(n+1)}}
=\frac{\sigma^2}{m^{n+2}}.
\]
Ortogonalność przyrostów daje
\[
\E W_n^2=1+\frac{\sigma^2}{m(m-1)}(1-m^{-n})
\leq1+\frac{\sigma^2}{m(m-1)}.
\]
Wniosek o zbieżności w $L^2$ i p.n. uzyskamy
w następnym rozdziale. Dla $m\leq1$ oszacowanie
nie daje jednostajnej kontroli drugiego momentu;
rozstrzygnięcie granic wymaga innego argumentu.
\end{przyklad}


